\documentclass[../../../include/open-logic-section]{subfiles}

\begin{document}

\olfileid{sth}{choice}{banach}
\olsection{د باناخ--تارسکي پاراډوکس}

د انتخاب د توجيې دپاره، لکه څنګه چې بولوس د بديلولو د
توجيې هڅه کړې وه، کېدے شي \emph{بهرنيو} دلايلو ته هم مخه کړو
(وګورئ \olref[replacement][extrinsic]{sec})۔ بالاخره د انتخاب منل
ډېرې خوښې پايلې لري: د هر دوو اصلي عددونو د پرتله کولو وړتيا؛
د هر سټ د ښه ترتيبولو وړتيا؛ د اصلي عددونو د ترتيبي عددونو د
يوه ځانګړي ډول په توګه ګڼلو وړتيا؛ او داسې نور۔

خو کله کله ويل کېږي چې انتخاب \emph{ناخوښې} پايلې هم لري۔
دا خبره تر ډېره د \cite{BanachTarski1924} د يوې پايلې له امله کېږي۔

\begin{thm}[د باناخ--تارسکي پاراډوکس (په $\ZFC$ کښې)]
هر توپ په متناهي شمېر ټوټو وېشل کېدے شي؛ بيا دا ټوټې
(په ګرځولو او لېږدولو سره) داسې بيا يوځاے کېدے شي چې
د هماغه توپ دوه نقلونه ترې جوړ شي۔
\end{thm}
\noindent
په لومړي نظر دا لږه حيرانوونکې ده۔ ښکاره ده چې دوه توپونه
د اصلي توپ د حجم \emph{دوه چنده} حجم لري۔ خو صُلب حرکتونه ---
ګرځول او لېږدول --- حجم نه بدلوي۔ نو داسې برېښي لکه
باناخ--تارسکي چې موږ ته له هېڅه نوې ماده راپيدا کوي۔

تر دې هم عجيبه خبره شته۔\footnote{وګورئ \citet[Theorem 3.12]{Wagon2016}۔}
ورته استدلال ښيي چې يوه وړه د نخود دانه په متناهي شمېر ټوټو
پرې کېدے شي، او بيا همدغه ټوټې (په ګرځولو او لېږدولو سره)
د بېګ بېن د شکل او اندازې يو څيز جوړولے شي۔

خو له دې څخه هېڅ يوه خبره يوازې په $\ZF$ کښې نه ثابتېږي۔\footnote{که څه هم
باناخ--تارسکي د انتخاب له اصله په ټينګه کمزورو اصلونو سره
ثابتېږي؛ وګورئ \citet[303]{Wagon2016}۔} نو له يوې پرېکړې سره
مخامخ يو: يا انتخاب رد کړو، يا له دې «پاراډوکس» سره ژوند ومنو۔

موږ وايو چې له دې «پاراډوکس» سره ژوند منل په کار دي۔ په
حقيقت کښې موږ ته دا ډېر پاراډوکس نه ښکاري۔ په ځانګړې توګه
نه وينو چې دا د لاندې پايلو په پرتله ولې زيات يا کم پاراډوکس
وګڼل شي:\footnote{\citet[276--7]{Potter2004}،
\citet[16]{Weston2003}، او \citet[31, 308--9]{Wagon2016}
له نورو بېلګو سره ورته ټکي بيانوي۔}
\begin{enumerate}
	\item په $(0,1)$ وقفه کښې د ټکو شمېر د $\Real$ د ټکو هومره دے۔
	\\\emph{ثبوت}: $\tan(\pi(r-\nicefrac{1}{2})))$ په پام کښې ونيسئ۔
	\item په يوه کرښه کښې د ټکو شمېر د يوه مربع د ټکو هومره دے۔
	\\وګورئ \olref[his][set][pathology]{sec} او \olref[his][set][cantorplane]{sec}۔
	\item فضا ډکوونکې منحنيانې شته۔
	\\وګورئ \olref[his][set][pathology]{sec} او \olref[his][set][hilbertcurve]{sec}۔
\end{enumerate}
له دې درېو پايلو څخه يوه هم انتخاب ته اړتيا نه لري۔ اوس
يې موږ د رياضياتو حيرانوونکې، ښکلې برخې ګڼو۔ ښايي د
باناخ--تارسکي پاراډوکس په اړه هم همداسې نظر ولرو۔
د چاپ شوې ټانجنټ تابعې د څرګندونې په پای کښې يو
زيات تړونکے قوس ښکاري؛ اصلي فورمول هماغسې ساتل شوے دے
(OLSTH-028)۔

البته دلته يوه تخنيکي خبره يادول پکار دي، خو په ارامه
ذهن پوهېدل پرې بس دي۔ صُلب حرکتونه حجم ساتي۔ نو هغه
پنځه\footnote{موږ پاراډوکس د «متناهي شمېر ټوټو» په ژبه بيان کړ۔
په حقيقت کښې \citet{Robinson1947} ثابته کړه چې وېش په
\emph{پنځو} ټوټو کېدے شي (خو تر پنځو په کمو نه)۔ د ثبوت
دپاره \citet[pp.~66--7]{Wagon2016} وګورئ۔} ټوټې چې توپ
ورباندې وېشل کېږي، ټولې \emph{د اندازې وړ} نه شي کېدے۔
نو په لنډو ټکو، دې هرې ټوټې ته په جلا توګه حجم ټاکل
معنا نه لري۔ دا ټوټې د ټکو داسې ناليدل کېدونکې «نامتناهي
خورې ورې ټولګې» وګڼئ۔ کېدے شي د دغسې «نامتناهي
خورو ورو» سټونو تصور عجيب وبرېښي۔ خو د هغو شتون له
هغه غوښتنې څخه راځي چې په \stagesacc{} کښې راغلې: د
مخکښې موجودو سټونو \emph{ټولې ممکنې} ټولګې جوړې کړئ۔

که بيا هم دا خبرې تاسو نه قانع کوي، د باناخ--تارسکي
پاراډوکس د منلو په ګټه دا وروستے (بهرنے) دليل واورئ۔ له دې
سمدستي د رياضياتو تر ټولو ښه ټوکه جوړېږي:
\begin{enumerate}
	\item[] \emph{پوښتنه}. د «باناخ--تارسکي» د تورو بل ترتيب څه دے؟
	\item[] \emph{ځواب}. «باناخ--تارسکي باناخ--تارسکي»۔
\end{enumerate}

\end{document}